
When a static analysis tool flags every failure in a test set, one might expect it to provide effective repair guidance. However, if those flags are overwhelmingly style conventions --- properties of LLM code generation context that fire on every generated snippet independent of whether the code is correct --- the LLM receiving such feedback is directed to address formatting while underlying logical errors remain undiagnosed. The result, as this study shows empirically, is not neutral: on HumanEval's algorithmic problems, pylint/mypy-guided repair actively reduces pass@1 by 4.3 percentage points relative to no-feedback single-pass generation.

This counterintuitive finding motivates the study. Iterative code repair --- generating code, receiving feedback on failures, and re-generating to correct them --- has emerged as an inference-time approach to improving LLM code quality without additional training [Chen et al., 2023; Shinn et al., 2023; Arimbur, 2026]. The practical question for system designers is which feedback signal to use: execution test results (run the code, observe failure messages), or static analysis (run pylint/mypy, observe warnings). Each approach has been studied in isolation, but no prior work provides a compute-controlled, head-to-head comparison of execution versus pylint/mypy feedback on functional correctness benchmarks, accompanied by mechanism analysis of why one signal outperforms the other.

The surface problem is well-established: LLMs generate incorrect code at substantial rates (39\% failure rate on HumanEval, 67\% on MBPP for a competitive 8B instruction-tuned model), and feedback-based repair can help reduce these failure rates. The deeper problem is that feedback signals differ in \textit{functional informativeness} --- the degree to which the signal diagnoses the specific errors present in a failing solution. Static analysis tools like pylint were designed to enforce code style conventions and detect syntactically suspicious patterns; they were not designed to diagnose the logical and runtime errors that dominate benchmark failures. No prior study has measured this informativeness difference empirically under compute-controlled conditions, nor decomposed which pylint signals are informative versus noise for functional correctness repair.

The key finding of this study is that pylint and mypy achieve high total \textit{coverage} of LLM code failures --- flagging every failing program --- but that the dominant signals are Convention-category style rules that fire regardless of functional correctness. Only 12.5\% of baseline failures receive a functional Error or Warning flag from pylint; mypy coverage is 0\%. This style-function dissociation explains both why pylint feedback fails to match execution feedback overall, and why it actively harms performance on complex algorithmic tasks: the LLM, acting on style guidance, may restructure solutions while introducing new logical errors.

This paper makes three contributions:

\textbf{C1: First iso-compute comparison of execution versus pylint/mypy feedback.} Execution test feedback is compared against pylint/mypy static analysis feedback at a fixed token budget of B=1000 output tokens per problem on HumanEval and MBPP using Llama 3.1 8B Instruct. Paired McNemar tests confirm execution feedback is significantly superior on both benchmarks (HumanEval: p=0.0001, MBPP: p<$10^{-18})$, with 15 problems uniquely repaired by execution versus zero by pylint on HumanEval, and 85 versus 2 on MBPP.

\textbf{C2: Style-function dissociation measurement.} Pylint flags are decomposed by category (E/W/C/R/I) across 64 HumanEval baseline failures, showing that 94.3\% are Convention-category style flags (C0304, C0114) that fire regardless of functional correctness, while only 12.5\% of failures receive a functional Error or Warning flag. Mypy coverage is 0\%. This decomposition reveals that coverage without category analysis systematically overstates feedback informativeness.

\textbf{C3: Benchmark asymmetry as evidence of task-complexity moderation.} Pylint/mypy repair harms HumanEval ($\Delta =-4.3pp)$ but benefits MBPP ($\Delta =+18.3pp)$, while execution feedback improves both. This asymmetry is interpreted as evidence that task complexity moderates the utility of style-guided repair: simple function-completion tasks tolerate style-triggered rewrites without losing logical structure, whereas complex algorithmic problems are sensitive to rewriting driven by style feedback.

The paper proceeds as follows. Section 2 reviews related work and positions the contribution. Section 3 describes the iso-compute methodology. Section 4 presents the experimental setup. Section 5 reports results. Section 6 discusses implications and limitations. Section 7 concludes.

